\chapter{Biopsy}
\index{Biopsy}

\section*{Definition and Overview}
A biopsy is a procedure in which a small sample of tissue is removed from the body so that it can be examined under a microscope by a pathologist. It is the most reliable way to determine whether an abnormal growth is cancer, and it is also used to diagnose many other conditions, including inflammation, infection, scarring, and transplant rejection.

Biopsies can be taken from almost any organ and are performed in a variety of ways, from a thin-needle sample to removal of an entire lesion. The pathology report generated from the sample provides diagnostic certainty and, increasingly, guides treatment decisions through the molecular testing of tumour tissue.

\section*{Epidemiology}
Biopsy is one of the most common diagnostic procedures in modern medicine, and its use continues to grow as screening programmes detect more abnormalities requiring characterisation.

\begin{itemize}
    \item \textbf{Volume:} many millions of biopsies are performed worldwide each year
    \item \textbf{Common indications:} a lump or mass, an abnormal finding on imaging such as a mammogram or CT scan, or a suspicious skin lesion
    \item \textbf{Safety:} the procedure is safe, and serious complications are uncommon
    \item \textbf{Evolving role:} biopsy tissue is increasingly used for molecular and genetic testing that guides targeted therapy
\end{itemize}

\section*{Aetiology and Pathophysiology}
A biopsy is performed when tissue must be examined directly to answer a diagnostic question that imaging and blood tests cannot resolve. The sample is fixed, embedded, sliced very thinly, and stained so that the cellular architecture can be seen clearly.

\begin{itemize}
    \item \textbf{Cancer diagnosis:} the most common reason, to confirm whether cells are malignant and to identify the tumour type, grade, and molecular features
    \item \textbf{Benign conditions:} to diagnose non-cancerous growths, inflammatory diseases, infection, or fibrosis
    \item \textbf{Organ disease:} for example kidney or liver biopsy to assess the severity of damage or the presence of rejection after transplantation
    \item \textbf{Sampling principle:} the sample must be representative of the abnormal area, so imaging such as ultrasound or CT is often used to guide the needle precisely
\end{itemize}

\section*{Clinical Features}
The features surrounding a biopsy relate both to the reason it is performed and to the procedure itself.

\begin{itemize}
    \item The test is often triggered by a new or growing lump, an abnormal scan, blood in the stool or urine, or a changing skin mole
    \item Most biopsies are performed as day procedures under local anaesthetic and take only minutes
    \item Mild pain, bruising, and soreness at the site are common afterwards and usually settle within a few days
    \item Some types require preparation, such as fasting, bowel preparation, or temporarily pausing blood-thinning medicines
\end{itemize}

\section*{Differential Diagnosis}
When a biopsy result is unclear, or when a lump or abnormal scan needs assessment, the possibilities that must be distinguished include:

\begin{itemize}
    \item \textbf{Benign versus malignant tumours:} the central question that most biopsies are designed to answer
    \item \textbf{Inadequate sample:} the needle may have missed the abnormality, requiring a repeat or larger biopsy
    \item \textbf{Borderline or pre-cancerous changes:} lesions that are not yet cancer but carry the potential to become so
    \item \textbf{Infection versus inflammation versus tumour:} each can produce a similar lump, discharge, or abnormal scan
    \item \textbf{Sampling error:} a false-negative result, in which the biopsy appears normal but the abnormality lies in adjacent tissue
\end{itemize}

\section*{When to Seek Medical Care}
See a doctor before a biopsy to discuss the procedure, benefits, and risks. After the biopsy, contact the clinic if you develop:

\begin{itemize}
    \item Heavy bleeding from the biopsy site that does not stop with pressure
    \item Increasing pain, redness, swelling, or discharge, which may indicate infection
    \item Fever or a general feeling of being unwell
    \item A wound that is not healing as expected
\end{itemize}

\textbf{Contact your local emergency services for:}
\begin{itemize}
    \item Chest pain or sudden difficulty breathing after a lung or chest biopsy, which may indicate a collapsed lung
    \item Collapse, dizziness, or fainting that may be due to bleeding
    \item Heavy, uncontrolled bleeding
    \item Swelling of the face or throat with difficulty breathing, suggesting a severe allergic reaction
\end{itemize}

\section*{Diagnosis and Investigations}
\begin{itemize}
    \item \textbf{Imaging-guided biopsy:} ultrasound, CT, or MRI is used to direct the needle precisely for lesions of the breast, liver, kidney, lung, and other deep organs
    \item \textbf{Pre-procedure checks:} blood tests for clotting, and a review of medicines such as anticoagulants that may need to be paused
    \item \textbf{Informed consent:} discussion of why the biopsy is needed, how it will be performed, and the associated risks
    \item \textbf{Methods of biopsy:} the technique depends on the site and the amount of tissue required:
    \begin{itemize}
        \item \textbf{Needle biopsy:} a thin needle collects cells (fine needle aspiration) or a small core of tissue (core needle biopsy)
        \item \textbf{Punch or shave biopsy:} for skin lesions
        \item \textbf{Endoscopic biopsy:} a sample taken through an endoscope, for example during colonoscopy or gastroscopy
        \item \textbf{Excisional biopsy:} removal of the entire lesion, often used for small lumps
        \item \textbf{Bone marrow biopsy:} a sample of bone marrow for the diagnosis of blood disorders
        \item \textbf{Sentinel lymph node biopsy:} removal of the first lymph node draining a cancer to check for spread
    \end{itemize}
    \item \textbf{Pathological examination:} the sample is processed and examined under a microscope, with special stains and molecular or genetic tests applied as needed
\end{itemize}

\section*{Management}
\begin{itemize}
    \item \textbf{The procedure itself:} most biopsies are performed under local anaesthetic; core and excisional biopsies may require a stitch and a small dressing
    \item \textbf{Post-procedure care:} keep the site clean and dry, apply firm pressure if it bleeds, and use simple pain relief as needed
    \item \textbf{Discussing results:} the treating doctor explains the pathology report and what it means for diagnosis and prognosis
    \item \textbf{Planning treatment:} the biopsy result determines whether further surgery, chemotherapy, radiotherapy, or active monitoring is appropriate
    \item \textbf{Repeat or additional biopsy:} arranged if the sample was inadequate, the result was uncertain, or more tissue is needed for molecular testing
    \item \textbf{Follow-up:} review of the biopsy site for healing and of the underlying condition
\end{itemize}

\section*{Complications}
\begin{itemize}
    \item Pain, bruising, and localised bleeding at the biopsy site
    \item Wound infection, which is uncommon and treated with antibiotics
    \item A collapsed lung (pneumothorax) after a needle biopsy of the lung, usually managed with observation or a chest drain
    \item Bleeding into an organ after a deep organ biopsy, which is rare but can be serious
    \item Damage to nearby structures or nerves at the biopsy site (uncommon)
    \item A false-negative result, in which the sample misses the abnormality and the diagnosis is delayed
\end{itemize}

\section*{Prognosis and Outcomes}
Biopsy is a safe and routine procedure, and for the great majority of people the outcome is a clear answer.

\begin{itemize}
    \item Most biopsies are completed within minutes, and recovery takes a day or two
    \item Serious complications are rare, and most minor ones settle without treatment
    \item The main value of biopsy lies in diagnostic certainty, which allows the appropriate treatment to be planned
    \item The overall outcome depends on the condition being investigated rather than on the biopsy itself
\end{itemize}

\section*{Prevention and Self-Care}
\begin{itemize}
    \item Follow all pre-procedure instructions, such as pausing blood thinners, fasting, or stopping certain medicines as advised
    \item Tell the doctor about allergies, all current medicines, and whether you might be pregnant
    \item After the procedure, keep the wound clean and dry and watch for signs of infection
    \item Attend the appointment to collect and discuss your results, and do not assume a benign result needs no follow-up
    \item Report any new symptoms, such as fever, chest pain, or heavy bleeding, promptly
\end{itemize}

\section*{Sources and Further Reading}
The following organisations provide reliable information on biopsy procedures for patients, families, and health professionals.

\begin{itemize}
    \item NHS. \textit{Information on biopsy procedures and what to expect.} https://www.nhs.uk/
    \item Mayo Clinic. \textit{Overview of biopsy types and how the procedure is performed.} https://www.mayoclinic.org/
    \item MedlinePlus. \textit{Health information on biopsy procedures for patients.} https://medlineplus.gov/
    \item National Cancer Institute. \textit{Educational resources on biopsy and cancer diagnosis.} https://www.cancer.gov/
    \item MSD Manual. \textit{Professional information on biopsy techniques and tissue diagnosis.} https://www.msdmanuals.com/
\end{itemize}
